% =====================================================================
% Seção 6 — Conclusão
% =====================================================================
\section{Conclusão}\label{sec:conc}

Retornamos à pergunta central: \textbf{existe um fator limítrofe entre produzir
IA com IA?} A evidência reunida --- teórica, estatística e empírica --- sustenta
uma resposta afirmativa e, ao mesmo tempo, matizada.

\textbf{Sim, existe um fator limítrofe estrutural.} Ele tem três manifestações
documentadas: (i) \emph{informacional} --- o treinamento recursivo sobre dados
gerados por modelo colapsa a distribuição aprendida, eliminando o suporte de
cauda longa e reduzindo a diversidade, fenômeno estatisticamente previsto e
confirmado empiricamente (\emph{Nature} 2024); (ii) \emph{de avaliação} --- quando
o mesmo processo gera e julga, o sinal de melhoria se corrompe
(\emph{self-play} hacking, discovery--reliability gap); e (iii) \emph{de âncora} ---
o progresso sustentado exige dados reais ou validação externa a cada iteração.

\textbf{Mas não é um muro absoluto; é um gargalo de informação e de feedback.} O
fator limítrofe não separa ``produzir IA com IA'' de ``não poder em absoluto''.
Separa a produção em \emph{loop fechado auto-suficiente} (que colapsa) da
produção \emph{com âncora externa} (que progride). O gargalo não está na
inteligência, mas na fonte de informação e na honestidade do sinal de avaliação.

O estudo de caso do \emph{OpenCode Ecosystem Core} (R459/R460) corporifica essa
tese em dados primários auditáveis: o HABD maximiza diversidade e cobertura
($\Div = 0{,}8333$, $\cob = 1{,}000$) mas viola o critério de relevância
($\Delta g/g = 0{,}104 > 0{,}05$), com veredito honesto \texttt{refuta\_H3}. O
resultado demonstra que o ecossistema --- que produz IA com IA continuamente ---
só o faz de forma sustentável porque cada loop é \emph{ancorado} por
especificação formal, testes externos ao gerador e auditoria independente.

\paragraph{Resposta direta à pergunta.}
Sim. Existe um fator limítrofe entre produzir IA com IA. Ele é real, tem
fundamento matemático e empírico, e é \emph{mitigável, não anulável}: a IA pode
produzir IA, desde que não seja \emph{apenas} IA produzindo IA. Sem âncora
externa de informação e de validação, o loop se retroalimenta até esgotar; com
âncora, ele progride. \emph{O limite não é a inteligência; é a fonte e o
verificador.}

\paragraph{Limites e anti-overclaim.}
Declaramos o escopo: as observações empíricas vêm de corpus-piloto controlado e
não generalizam para produção. Distinguimos ``ganho observável'' (fato) de
``satisfaz critério rígido de relevância'' (não satisfeito). Não afirmamos que o
fator limítrofe seja inviolável para sempre, nem que os resultados de
distribuições simuladas se transfiram integralmente a linguagem natural. O
artigo oferece, portanto, uma resposta fundamentada, equilibrada e honesta ---
consistente com o princípio de \emph{anti-overclaim} que governa o ecossistema
que o produziu.
